\documentclass[11pt]{article}
% Shared packages for SI table/figure masters (input by each master).
\usepackage[margin=0.6in]{geometry}
\usepackage{booktabs}
\usepackage{longtable}
\usepackage{caption}
\usepackage{hyperref}
\usepackage{graphicx}
\usepackage{float}
\usepackage{amsmath}
\usepackage{enumitem}
\graphicspath{{../../figures/}{../}{./}}
\setlength{\parindent}{0pt}
\pagestyle{plain}

\setlength{\parskip}{0.85em}
\begin{document}

\begin{center}
{\LARGE\bfseries Supplementary Text}
\end{center}

\vspace{0.4em}
\noindent\textbf{The CSPdb}

The CSPdb includes receptors that regulate diverse biological processes, including
cell cycle progression (e.g., MAD2: 1KLQ), transcriptional control (e.g., Nrd1:
2MOW, 6GC3), translational regulation, autophagy, apoptosis (e.g., BCL-XL family:
2M04, 2LP8), and cell adhesion (e.g., tensin: 2LOZ).
Disease-relevant targets address mechanisms in leukemia and other cancers,
Alzheimer's disease (e.g., amyloid precursor protein-binding
proteins: 2ROZ), and infectious diseases from viruses and bacteria, such as
anthrax (2RUI), chikungunya (5I22), SARS-CoV-2 (7PKU), HIV (2L6E), and Ebola
(2KQ0).

The CSPdb features several recurring protein domains that exhibit notable
conformational changes upon ligand binding. For instance, it includes 23
interactions involving calmodulin (CaM) or CaM-like proteins with peptides
(e.g., PDB IDs: 1SY9, 2JZI; full list in \textbf{SI Table S1}), which often involve
substantial rearrangements upon binding described in detail
elsewhere\textsuperscript{1}. Other receptors known for binding-induced
conformational shifts or folding include the G protein Cdc42 (1CF4), the
metaphase chromosome spindle checkpoint protein MAD2 (1KLQ), the KIX domain
(2LXS), and CREB-binding proteins (e.g., 1R8U, 1L8C).

Additional recurring domains include the extra-terminal (ET) domains of
bromodomain proteins BRD3 and BRD4 (e.g., 2NCZ, 2ND0; full list in \textbf{SI Table S7},
\textbf{SI Fig.\ S6}), which bind intrinsically-disordered regions of viral integrase or
chromatin regulator proteins. These interactions trigger subtle
changes distal from the binding site, including the formation of a stabilizing
secondary structure that reorients two allosteric helices\textsuperscript{2, 3}.
Similarly, the PDZ2 domain (e.g., 1D5G, 1VJ6, 2M0V) shows distal effects
mediated by the $\beta$2--$\beta$3 loop during peptide
recognition\textsuperscript{4}.

Nine entries involve the PH domain of the TFIIH transcription initiation factor
subunit Tfb1 (e.g., 2LOX, 2M14; full list in \textbf{SI Table S8}, \textbf{SI Fig.\ S7}). In these
complexes, the ligand is often modeled as a helical polypeptide bound
orthogonally to an extended $\beta$-sheet, resulting in signal
propagation via hydrogen-bond rearrangements. Other targets which involve
hydrogen bond rearrangement in $\beta$-sheets include the Cdc42 (1CF4),
Cyclophilin A (2MS4), and SH3 domain (e.g., 2OJ2, 2ROL; full list in \textbf{SI Table S1}). An additional eight entries use ubiquitin or ubiquitin-like (e.g.\ SUMO
and NEDD8) protein receptors. Ubiquitin has been used as a model system for NMR
studies\textsuperscript{5} and for studying allostery\textsuperscript{6}
(e.g., 2KWV, 2MUR; full list in \textbf{SI Table S9}, \textbf{SI Fig.\ S8}).

\vspace{0.4em}
\noindent\textbf{Estimated Distance Limits for Non-structural NMR Chemical Shift Perturbations}

Aromatic ring-current effects decrease as $r^{-3}$, with established
ring-current parameters indicating that a $\geq 0.05$-ppm amide $^{1}$H shift is
unlikely beyond approximately 10\,\AA{} for Phe/Tyr and 12\,\AA{} for Trp, even
under idealized geometry\textsuperscript{8, 9}. Similarly, changes in charge or
ionization state accompanying ligand binding can produce amide CSPs through
changes in the local electric field, although the magnitude of these effects
depends strongly on the dielectric environment, screening, and
orientation\textsuperscript{10, 11}. Thus, distal CSPs (e.g., $>15$\,\AA{} from
the binding site) generally reflect long-range effects of ligand binding and are
consistent with changes in protein conformation or conformational distributions.

\vspace{0.4em}
\noindent\textbf{NMR ring-current contributions to chemical shift perturbations}

Based on the $r^{-3}$ distance dependence of aromatic ring-current effects and
established ring-current parameters for aromatic amino acids, the theoretical
maximum distance for a $\geq 0.05$-ppm HN ring-current shift is approximately
10\,\AA{} for Phe/Tyr and 12\,\AA{} for Trp under idealized
geometry\textsuperscript{8, 9}.

\textbf{Estimated maximum distance for a $\geq 0.05$-ppm amide $^{1}$H ring-current shift.}

Aromatic ring-current contributions to $^{1}$H chemical shifts can be
approximated using the equivalent-dipole model,
\begin{equation*}
\Delta\delta_{\mathrm{RC}}=\mu\frac{(3\cos^{2}\theta-1)}{r^{3}},
\end{equation*}
where $r$ is the distance from the proton to the aromatic ring center,
$\theta$ is the angle between this vector and the normal to the ring plane, and
$\mu$ is the effective ring-current constant. Abraham et al.\ determined
$\mu = 26.2$\,ppm\,\AA{}$^{3}$ for benzene. The maximum absolute effect occurs
at $\theta=0^{\circ}$ or $180^{\circ}$ corresponding to the proton lying along
the normal above or below the aromatic ring plane, where the angular term has a
magnitude of 2. Using the 0.05-ppm significance threshold gives
\begin{equation*}
r_{\max}=\left(\frac{2\mu}{0.05}\right)^{1/3}=10.2\,\textrm{\AA}.
\end{equation*}
Thus, even under idealized geometry, a Phe/Tyr-like aromatic ring would need to
be within approximately 10\,\AA{} of an amide proton to produce a significant
ring-current effect. For Trp, the contributions from its two aromatic rings are
summed; using the relative ring-current intensities reported by Perkins et al.\
gives an idealized upper-bound estimate of approximately 12\,\AA{}.

\vspace{0.4em}
\noindent\textbf{NMR electric-field induced chemical shift perturbations}

Based on the $r^{-2}$ distance dependence of electric-field effects and
experimentally determined parameters for protein amide chemical shifts, a
one-unit charge change is estimated to produce a $\geq 0.05$-ppm HN shift only
within approximately 6--8\,\AA{} under idealized
conditions\textsuperscript{10, 11}.

\textbf{Estimated maximum distance for a $\geq 0.05$-ppm amide $^{1}$H electric-field induced shift.}

Electric-field-induced amide chemical shifts can be approximated using the
Buckingham treatment,
\begin{equation*}
|\Delta\delta_{\mathrm{HN}}| = |A|\,|E|\,|\cos\beta|,
\end{equation*}
where $A$ is the HN shielding polarizability, $E$ is the change in electric
field at the amide proton, and $\beta$ is the angle between these vectors.
Using $|A|=21$\,ppm\,\AA{}$^{2}/e$ for backbone amide protons\textsuperscript{11}
and approximating the electric field from a point-charge change using Coulomb's
law gives
\begin{equation*}
|\Delta\delta_{\mathrm{HN}}| = \frac{21\,|\Delta q|}{\epsilon_{\mathrm{eff}}\,r^{2}}\,|\cos\beta|.
\end{equation*}
For a one-unit charge change ($|\Delta q|=1e$), optimal orientation
($|\cos\beta|=1$), and the 0.05-ppm significance threshold,
\begin{equation*}
r_{\max}=\left[\frac{21}{0.05\,\epsilon_{\mathrm{eff}}}\right]^{1/2}.
\end{equation*}
Experimentally derived effective dielectric constants for proteins are
approximately 6--13; values of $\sim$6.5--7 and $8.6\pm 0.8$ have been obtained
from protein CSP measurements\textsuperscript{10, 11}. These values correspond
to idealized maximum distances of approximately 6--8\,\AA{} for a
$\geq 0.05$-ppm $^{1}$HN shift from a one-unit charge change.

Unlike aromatic ring-current effects, a universal distance cutoff cannot be
assigned to electrostatic CSPs because their magnitude depends on charge,
orientation, dielectric environment, and screening. Small electric-field-induced
$^{1}$HN CSPs have been experimentally detected $>12$\,\AA{} from a charge
perturbation, although these were below the 0.05-ppm significance threshold used
here\textsuperscript{11}. Although our apo and holo datasets were matched for
pH, ligand binding itself may alter protein ionization states and thereby
generate electric-field-induced CSPs.

\end{document}
